\ifdefined\standalone 
\documentclass[12pt]{article}
\usepackage{graphicx}
\usepackage{float}
\usepackage[a4paper, margin=1in]{geometry}
\usepackage[utf8]{inputenc}
\usepackage{textcomp}
\usepackage{amsmath}
\usepackage{booktabs}
\usepackage{tikz}
\usepackage{enumitem}
\usetikzlibrary{decorations.pathmorphing,patterns}
\usepackage[hidelinks]{hyperref} % <- hypertext

\begin{document}
\fi

\section{Mixture law}
\label{sec:mixture-law}

This verification case evaluates the ability of \textsc{Gpyro} to apply the mixture law for reacting condensed-phase materials. For a generic property $\phi$ (such as bulk density $\rho$, porosity $\Psi$, thermal conductivity $k$ or the emissivity $\varepsilon$), the effective mixture value is written as:
\begin{equation}
\bar{\phi}(x,t) = \sum_{i=1}^{M} X_i(x,t)\, \phi_i ,
\end{equation}
where $X_i$ is the volume fraction of material $i$.


A simulation is performed in which material A is converted into material B at a constant reaction rate
$r = 0.001~\mathrm{kg\,m^{-3}\,s^{-1}}$. The material properties are summarized in Table~\ref{tab:properties_mixture_law}.
The analytical solution for the mass and volume fractions $Y$ is given by:
\begin{equation}
\left\{
\begin{aligned}
Y_A(t) &= \max(1 - r t,\, 0), \\
Y_B(t) &= 1 - Y_A(t),
\end{aligned}
\right.
\qquad
\left\{
\begin{aligned}
X_A(t) &= \bar{\rho}\, {Y_A}/{\rho_A}, \\
X_B(t) &= \bar{\rho}\, {Y_B}/{\rho_B},
\end{aligned}
\right
\qquad
\text{with }
\bar{\rho}(t) =
\left( \frac{Y_A}{\rho_A} + \frac{Y_B}{\rho_B} \right)^{-1}
\end{equation}
\begin{table}[H]
\centering
\caption{Thermophysical properties used for the mixture law verification}
\label{tab:properties_mixture_law}
\renewcommand{\arraystretch}{1.3}
\begin{tabular}{c c c c c c}
\hline
Material &
$k$ ($\mathrm{W\,m^{-1}\,K^{-1}}$) &
$\rho$ ($\mathrm{kg\,m^{-3}}$) &
$\rho_s$ ($\mathrm{kg\,m^{-3}}$) &
$c_p$ ($\mathrm{J\,kg^{-1}\,K^{-1}}$) &
$\varepsilon$ ($-$) \\
\hline
Material A & 1.0 & 1000 & 2000 & 1200 & 1.0 \\
Material B & 0.5 & 600  & 800  & 1900 & 0.0 \\
\hline
\end{tabular}
\end{table}

The temporal evolution of mass fractions, volume fractions, and effective properties is shown in Figure~\ref{fig:Mix_law_K}. Numerical results obtained with the current version of \textsc{Gpyro} are compared to the analytical solution and to results from \textsc{Gpyro} V0.8, highlighting the correct implementation of the mixture law in the version of \textsc{Gpyro}.

\begin{figure}[H]
\centering
\includegraphics[width=1\linewidth]{All_Variables_Subplots_ML.png}
\caption{Mixture law verification: comparison between analytical and numerical results}
\label{fig:Mix_law_K}
\end{figure}

\ifdefined\standalone
\end{document}
\fi






